\chapter{INTRODUÇÃO}
\label{cap:introducao}

A área de ciência e engenharia dos materiais abrange metalurgia, semicondutores, engenharia cerâmica e ciência dos polímeros. É um campo multidisciplinar que viabiliza a existência de novas tecnologias necessárias para enfrentar uma grande variedade de desafios críticos da sociedade \cite{Thornton2009}.


\lipsum[1] %Apagar esta linha - Gerador de texto aleatório

Desenvolver uma contextualização consistente aqui.

\section{Justificativa}

Dado o contexto apresentado anteriormente, percebe-se que construir um \lipsum[2] %Apagar esta linha - Gerador de texto aleatório

\lipsum[3] %Apagar esta linha - Gerador de texto aleatório

\section{Objetivos}

\lipsum[4] %Apagar esta linha - Gerador de texto aleatório

\subsection{Objetivos Específicos}

Dentre os objetivos específicos deste trabalho pode-se destacar:
\begin{itemize}

\item Criar uma base de conhecimento de PADI aplicado à caracterização...; 

\item Desenvolver módulos para o...;

\item Implementar recursos para permitir que... ;

\item Criar dois assistentes de... ;

\item Gerar dois estudos de caso para validar os assistentes... ;

\item Viabilizar o uso do software por profissionais da área de... ;

\end{itemize}

\section{Estrutura do Trabalho}

Este trabalho está dividido em sete capítulos.  O \autoref{cap:introducao} expõe o contexto do estudo, as justificativas desta pesquisa e os objetivos a serem atingidos. O \autoref{cap:fundamentacao} apresenta conceitos de caracterização de materiais, relata algumas normas de organizações internacionais para este fim e faz uma explanação sobre a área de processamento e análise de imagens. O \autoref{cap:trabalhos}... . Finalmente, o \autoref{cap:conclusão} apresenta a conclusão e os potenciais trabalhos futuros a serem desenvolvidos.